% questions.tex —— 2013 年·福建卷·理科 · 题目内容片段（无答案）
% 由 tools/generate_projects.py 依来源总账自 2013/2013_数学真题_福建理_试题.md 生成，勿手改（改源 MD 后重跑）


\begin{gknotice}
\textbf{来源说明}：本题面依据固定源 PDF 的两页 A3 题面逐题转录，原卷印刷页码为 799—800。第 799 面为选择题 1—10、填空题 11—15、解答题 16，第 800 面为解答题 17—21。第 21 题为三选二（A 矩阵与变换、B 坐标系与参数方程、C 不等式选讲），三题全部收录。

\textbf{解析归属}：固定源不附答案或解析；配套答案与解析均由本库按本题面独立复算。

\textbf{图形说明}：第 4 题频率分布直方图、第 6 题程序框图、第 12 题三视图、第 13 题三角形、第 18 题正方形与抛物线构造、第 19 题四棱柱，均由项目内 TikZ 矢量源重绘。
\end{gknotice}

\gkpart{一、选择题}

\begin{question}
\gkti{1}已知复数 $z$ 的共轭复数 $\bar z=1+2\mathrm i$（$\mathrm i$ 为虚数单位），则 $z$ 在复平面内对应的点位于\mbox{（　　）}

（A）第一象限；（B）第二象限；（C）第三象限；（D）第四象限。

\end{question}
\begin{question}
\gkti{2}已知集合 $A=\{1,a\}$，$B=\{1,2,3\}$，则“$a=3$”是“$A\subseteq B$”的\mbox{（　　）}

（A）充分而不必要条件；（B）必要而不充分条件；（C）充分必要条件；（D）既不充分也不必要条件。

\end{question}
\begin{question}
\gkti{3}双曲线 $\dfrac{x^2}{4}-y^2=1$ 的顶点到渐近线的距离等于\mbox{（　　）}

（A）$\dfrac25$；（B）$\dfrac45$；（C）$\dfrac{2\sqrt5}{5}$；（D）$\dfrac{4\sqrt5}{5}$。

\Needspace{0.4\textheight}
\end{question}
\begin{question}
\gkti{4}某校从高一年级学生中随机抽取部分学生，将他们的模块测试成绩分成 $6$ 组：$[40,50)$，$[50,60)$，$[60,70)$，$[70,80)$，$[80,90)$，$[90,100)$ 加以统计，得到如图所示的频率分布直方图。已知高一年级共有学生 $600$ 名，据此估计，该模块测试成绩不少于 $60$ 分的学生人数为\mbox{（　　）}
\nopagebreak[4]
\begin{center}% ============================================================
% f2013_fujian_science_q04_histogram.tex
% 2013 福建卷（理工类）第 4 题 频率分布直方图（TikZ 矢量重绘）
% 组距 10，各组矩形高度按原页：0.005, 0.010, 0.015, 0.030, 0.025, 0.010
% 横轴 40--100 分；60 分以上四组高度之和为 0.080
% ============================================================
\begin{tikzpicture}[
  x=1cm, y=1cm, line width=0.6pt, font=\footnotesize, >=Stealth
]

  % ---------- 坐标轴 ----------
  \draw[->] (0,0) -- (8.9,0) node[below left, yshift=2pt] {分数};
  \draw[->] (0,0) -- (0,2.85) node[above] {};
  \node[align=center] at (0.58,3.10) {频率\\[-2pt]组距};
  \draw (0.18,3.08) -- (0.98,3.08);

  % 横轴断点符号
  \draw (0.32,0) -- (0.40,0.10) -- (0.48,0) -- (0.56,0.10) -- (0.64,0);

  % ---------- 纵向刻度与参考虚线 ----------
  \foreach \h/\lab in {0.005/0.005,0.010/0.010,0.015/0.015,0.025/0.025,0.030/0.030} {
    \pgfmathsetmacro{\yy}{\h*80}
    \draw[dashed, line width=0.35pt, gray!70] (0.08,\yy) -- (8.25,\yy);
    \node[left] at (-0.05,\yy) {\lab};
    \draw (-0.05,\yy) -- (0.05,\yy);
  }
  \node[left] at (-0.05,0) {$0$};

  % ---------- 各矩形 ----------
  % 40--50: 0.005
  \draw (3.20,0) rectangle (4.00,0.40);
  % 50--60: 0.010
  \draw (4.00,0) rectangle (4.80,0.80);
  % 60--70: 0.015
  \draw (4.80,0) rectangle (5.60,1.20);
  % 70--80: 0.030
  \draw (5.60,0) rectangle (6.40,2.40);
  % 80--90: 0.025
  \draw (6.40,0) rectangle (7.20,2.00);
  % 90--100: 0.010
  \draw (7.20,0) rectangle (8.00,0.80);

  % ---------- 横轴刻度 ----------
  \foreach \s/\xx in {40/3.20,50/4.00,60/4.80,70/5.60,80/6.40,90/7.20,100/8.00} {
    \draw (\xx,0) -- (\xx,-0.08);
    \node[below] at (\xx,-0.10) {\s};
  }

\end{tikzpicture}
\end{center}
\nopagebreak[4]
（A）$588$；（B）$480$；（C）$450$；（D）$120$。

\end{question}
\begin{question}
\gkti{5}满足 $a,b\in\{-1,0,1,2\}$，且关于 $x$ 的方程 $ax^2+2x+b=0$ 有实数解的有序数对 $(a,b)$ 的个数为\mbox{（　　）}

（A）$14$；（B）$13$；（C）$12$；（D）$10$。

\Needspace{0.4\textheight}
\end{question}
\begin{question}
\gkti{6}阅读如图所示的程序框图，若输入的 $k=10$，则该算法的功能是\mbox{（　　）}
\nopagebreak[4]
\begin{center}% ============================================================
% f2013_fujian_science_q06_flowchart.tex
% 2013 福建卷（理工类）第 6 题 程序框图（TikZ 矢量重绘）
% S=0, i=1；输入 k；S=1+2S；i=i+1；i>k？否则回代，是则输出 S
% ============================================================
\begin{tikzpicture}[
  x=1cm, y=1cm, line width=0.6pt, font=\footnotesize, >=Stealth,
  every node/.style={inner sep=3pt},
  proc/.style={rectangle, draw, minimum width=2.6cm, minimum height=0.62cm},
  io/.style={trapezium, draw, trapezium left angle=70, trapezium right angle=110,
             minimum width=2.8cm, minimum height=0.66cm},
  dec/.style={diamond, draw, aspect=2.2, minimum width=2.8cm, minimum height=0.85cm},
  term/.style={rectangle, rounded corners=8pt, draw, minimum width=2.0cm, minimum height=0.58cm}
]

  \node[term] (start) at (0,0) {开始};
  \node[proc] (init)  at (0,-1.35) {$S=0,\ i=1$};
  \node[io]   (input) at (0,-2.60) {输入 $k$};
  \node[proc] (upd)   at (0,-3.85) {$S=1+2S$};
  \node[proc] (inc)   at (0,-5.10) {$i=i+1$};
  \node[dec]  (judge) at (0,-6.55) {$i>k$?};
  \node[io]   (out)   at (0,-8.10) {输出 $S$};
  \node[term] (end)   at (0,-9.35) {结束};

  \draw[->] (start) -- (init);
  \draw[->] (init) -- (input);
  \draw[->] (input) -- (upd);
  \draw[->] (upd) -- (inc);
  \draw[->] (inc) -- (judge);
  \draw[->] (judge) -- node[right] {是} (out);
  \draw[->] (out) -- (end);

  % 否：回到 S=1+2S 之前
  \draw[->] (judge.east) -- (2.60,-6.55) -- (2.60,-3.85) -- (upd.east);
  \node[above] at (1.30,-6.45) {否};

\end{tikzpicture}
\end{center}
\nopagebreak[4]
（A）计算数列 $\{2^{n-1}\}$ 的前 $10$ 项和；（B）计算数列 $\{2^{n-1}\}$ 的前 $9$ 项和；（C）计算数列 $\{2^n-1\}$ 的前 $10$ 项和；（D）计算数列 $\{2^n-1\}$ 的前 $9$ 项和。

\end{question}
\begin{question}
\gkti{7}在四边形 $ABCD$ 中，$\overrightarrow{AC}=(1,2)$，$\overrightarrow{BD}=(-4,2)$，则该四边形的面积为\mbox{（　　）}

（A）$\sqrt5$；（B）$2\sqrt5$；（C）$5$；（D）$10$。

\end{question}
\begin{question}
\gkti{8}设函数 $f(x)$ 的定义域为 $\mathbf R$，$x_0$（$x_0\neq0$）是 $f(x)$ 的极大值点，以下结论一定正确的是\mbox{（　　）}

（A）$\forall x\in\mathbf R$，$f(x)\leqslant f(x_0)$；（B）$-x_0$ 是 $f(-x)$ 的极小值点；（C）$-x_0$ 是 $-f(x)$ 的极小值点；（D）$-x_0$ 是 $-f(-x)$ 的极小值点。

\end{question}
\begin{question}
\gkti{9}已知等比数列 $\{a_n\}$ 的公比为 $q$，记
\nopagebreak[4]
\[b_n=a_{m(n-1)+1}+a_{m(n-1)+2}+\cdots+a_{m(n-1)+m},\qquad c_n=a_{m(n-1)+1}\cdot a_{m(n-1)+2}\cdot\cdots\cdot a_{m(n-1)+m},\]
其中 $m,n\in\mathbf N^*$，则以下结论一定正确的是\mbox{（　　）}

（A）数列 $\{b_n\}$ 为等差数列，公差为 $q^m$；（B）数列 $\{b_n\}$ 为等比数列，公比为 $q^{2m}$；（C）数列 $\{c_n\}$ 为等比数列，公比为 $q^{m^2}$；（D）数列 $\{c_n\}$ 为等比数列，公比为 $q^{m^m}$。

\end{question}
\begin{question}
\gkti{10}设 $S$，$T$ 是 $\mathbf R$ 的两个非空子集，如果存在一个从 $S$ 到 $T$ 的函数 $y=f(x)$ 满足：

① $T=\{f(x)\mid x\in S\}$；

② 对任意 $x_1,x_2\in S$，当 $x_1<x_2$ 时，恒有 $f(x_1)<f(x_2)$；

那么称这两个集合“保序同构”。以下集合对不是“保序同构”的是\mbox{（　　）}

（A）$A=\mathbf N^*$，$B=\mathbf N$；（B）$A=\{x\mid-1\leqslant x\leqslant3\}$，$B=\{x\mid x=-8\text{ \text{或} }0<x\leqslant10\}$；（C）$A=\{x\mid0<x<1\}$，$B=\mathbf R$；（D）$A=\mathbf Z$，$B=\mathbf Q$。

\end{question}
\gkpart{二、填空题}

\begin{question}
\gkti{11}利用计算机产生 $0\sim1$ 之间的均匀随机数 $a$，则事件“$3a-1>0$”发生的概率为\gkblank{}。

\Needspace{0.4\textheight}
\end{question}
\begin{question}
\gkti{12}已知某一多面体内接于球构成一个简单组合体，如果该组合体的正视图、侧视图、俯视图均如图所示，且图中的四边形是边长为 $2$ 的正方形，则该球的表面积是\gkblank{}。
\nopagebreak[4]
\begin{center}% ============================================================
% f2013_fujian_science_q12_views.tex
% 2013 福建卷（理工类）第 12 题 三视图（TikZ 矢量重绘）
% 正视图、侧视图、俯视图均为：圆内含正方形（虚线），边长 2
% 正方体边长 2，外接球半径 sqrt(3)，表面积 12*pi
% ============================================================
\begin{tikzpicture}[
  x=1cm, y=1cm, line width=0.6pt, font=\footnotesize, >=Stealth
]

  \foreach \dx/\lab in {0/正视图, 4.05/侧视图, 8.10/俯视图} {
    \begin{scope}[xshift=\dx cm]
      % 外接圆
      \draw (0,0) circle (1.55);
      % 正方形（虚线）
      \draw[dashed] (-1,-1) rectangle (1,1);
      % 边长标注
      \draw[<->] (-1,-1.22) -- (1,-1.22);
      \node[below] at (0,-1.24) {$2$};
      % 视图名称
      \node at (0,-2.05) {\lab};
    \end{scope}
  }

\end{tikzpicture}
\end{center}
\nopagebreak[4]
\Needspace{0.4\textheight}
\end{question}
\begin{question}
\gkti{13}如图，在 $\triangle ABC$ 中，已知点 $D$ 在 $BC$ 边上，$AD\perp AC$，$\sin\angle BAC=\dfrac{2\sqrt2}{3}$，$AB=3\sqrt2$，$AD=3$，则 $BD$ 的长为\gkblank{}。
\nopagebreak[4]
\begin{center}% ============================================================
% f2013_fujian_science_q13_triangle.tex
% 2013 福建卷（理工类）第 13 题 三角形示意图（TikZ 矢量重绘）
% A 在上，D 在 BC 上，AD 垂直 AC
% 取 A(0,0), D(0,-3), C(3sqrt2,0), B(-sqrt2,-4)
% 整体旋转 -20 度以贴近原页版式
% ============================================================
\begin{tikzpicture}[
  x=1cm, y=1cm, line width=0.6pt, font=\footnotesize, >=Stealth
]

  \begin{scope}[rotate=-20]
    \coordinate (A) at (0,0);
    \coordinate (B) at (-1.4142,-4);
    \coordinate (C) at (4.2426,0);
    \coordinate (D) at (0,-3);

    \draw (A) -- (B) -- (C) -- cycle;
    \draw (A) -- (D);

    % 直角标记（AD 与 AC）
    \draw (0,-0.38) -- (0.38,-0.38) -- (0.38,0);

    \node[above] at (A) {$A$};
    \node[below left] at (B) {$B$};
    \node[below right] at (C) {$C$};
    \node[below left] at (D) {$D$};
  \end{scope}

\end{tikzpicture}
\end{center}
\nopagebreak[4]
\end{question}
\begin{question}
\gkti{14}椭圆 $\Gamma:\dfrac{x^2}{a^2}+\dfrac{y^2}{b^2}=1$（$a>b>0$）的左右焦点分别为 $F_1$，$F_2$，焦距为 $2c$。若直线 $y=\sqrt3(x+c)$ 与椭圆 $\Gamma$ 的一个交点 $M$ 满足 $\angle MF_1F_2=2\angle MF_2F_1$，则该椭圆的离心率等于\gkblank{}。

\end{question}
\begin{question}
\gkti{15}当 $x\in\mathbf R$，$|x|<1$ 时，有如下表达式：
\nopagebreak[4]
\[1+x+x^2+\cdots+x^n+\cdots=\frac1{1-x}.\]
\Needspace{5\baselineskip}
两边同时积分得：
\nopagebreak[4]
\[\int_0^{\frac12}1\,\mathrm dx+\int_0^{\frac12}x\,\mathrm dx+\int_0^{\frac12}x^2\,\mathrm dx+\cdots+\int_0^{\frac12}x^n\,\mathrm dx+\cdots=\int_0^{\frac12}\frac1{1-x}\,\mathrm dx,\]
\Needspace{5\baselineskip}
从而得到如下等式：
\nopagebreak[4]
\[1\times\frac12+\frac12\times\left(\frac12\right)^2+\frac13\times\left(\frac12\right)^3+\cdots+\frac1{n+1}\times\left(\frac12\right)^{n+1}+\cdots=\ln2.\]
\Needspace{5\baselineskip}
请根据以上材料所蕴含的数学思想方法，计算：
\nopagebreak[4]
\[\mathrm C_n^0\times\frac12+\frac12\mathrm C_n^1\times\left(\frac12\right)^2+\frac13\mathrm C_n^2\times\left(\frac12\right)^3+\cdots+\frac1{n+1}\mathrm C_n^n\times\left(\frac12\right)^{n+1}=\underline{\qquad\qquad}.\]

\end{question}
\gkpart{三、解答题}

\begin{question}
\gkti{16}某联欢晚会举行抽奖活动，举办方设置了甲、乙两种抽奖方案，方案甲的中奖率为 $\dfrac23$，中奖可以获得 $2$ 分；方案乙的中奖率为 $\dfrac25$，中奖可以获得 $3$ 分；未中奖则不得分。每人有且只有一次抽奖机会，每次抽奖中奖与否互不影响，晚会结束后凭分数兑换奖品。

（1） 若小明选择方案甲抽奖，小红选择方案乙抽奖，记他们的累计得分为 $X$，求 $X\leqslant3$ 的概率；

（2） 若小明、小红两人都选择方案甲或都选择方案乙进行抽奖，问：他们选择何种方案抽奖，累计得分的数学期望较大？

\end{question}
\begin{question}
\gkti{17}已知函数 $f(x)=x-a\ln x$（$a\in\mathbf R$）。

（1） 当 $a=2$ 时，求曲线 $y=f(x)$ 在点 $A(1,f(1))$ 处的切线方程；

（2） 求函数 $f(x)$ 的极值。

\Needspace{0.4\textheight}
\end{question}
\begin{question}
\gkti{18}如图，在正方形 $OABC$ 中，$O$ 为坐标原点，点 $A$ 的坐标为 $(10,0)$，点 $C$ 的坐标为 $(0,10)$。分别将线段 $OA$ 和 $AB$ 十等分，分点分别记为 $A_1,A_2,\cdots,A_9$ 和 $B_1,B_2,\cdots,B_9$，连接 $OB_i$，过 $A_i$ 作 $x$ 轴的垂线与 $OB_i$ 交于点 $P_i$（$i\in\mathbf N^*$，$1\leqslant i\leqslant9$）。

（1） 求证：点 $P_i$（$i\in\mathbf N^*$，$1\leqslant i\leqslant9$）都在同一条抛物线上，并求该抛物线 $E$ 的方程；

（2） 过点 $C$ 作直线 $l$ 与抛物线 $E$ 交于不同的两点 $M$，$N$，若 $\triangle OCM$ 与 $\triangle OCN$ 的面积之比为 $4:1$，求直线 $l$ 的方程。
\nopagebreak[4]
\begin{center}% ============================================================
% f2013_fujian_science_q18_parabola.tex
% 2013 福建卷（理工类）第 18 题 正方形与抛物线构造（TikZ 矢量重绘）
% O(0,0), A(10,0), B(10,10), C(0,10)
% OA、AB 各十等分；A_i(i,0)，B_i(10,i)；OB_i 与过 A_i 的竖线交于 P_i
% P_i(i, i^2/10) 在抛物线 y=x^2/10 上
% ============================================================
\begin{tikzpicture}[
  x=0.42cm, y=0.42cm, line width=0.6pt, font=\footnotesize, >=Stealth
]

  % ---------- 正方形 ----------
  \draw (0,0) -- (10,0) -- (10,10) -- (0,10) -- cycle;

  % ---------- 十等分刻度 ----------
  \foreach \i in {1,...,9} {
    \draw (\i,0) -- (\i,-0.22);
    \draw (10,\i) -- (10.22,\i);
  }

  % ---------- OB_i 与竖线 ----------
  \foreach \i in {1,...,9} {
    \draw[gray!65, line width=0.35pt] (0,0) -- (10,\i);
    \draw[dashed, gray!65, line width=0.35pt] (\i,0) -- (\i,{\i*\i/10});
    \fill (\i,{\i*\i/10}) circle (0.09);
  }

  % ---------- 标注 ----------
  \node[below left]  at (0,0)   {$O$};
  \node[below right] at (10,0)  {$A$};
  \node[above right] at (10,10) {$B$};
  \node[above left]  at (0,10)  {$C$};

  \node[below] at (1,-0.28) {$A_1$};
  \node[below] at (2,-0.28) {$A_2$};
  \node[below] at (5,-0.28) {$A_i$};
  \node[below] at (9,-0.28) {$A_9$};

  \node[right] at (10.28,1) {$B_1$};
  \node[right] at (10.28,2) {$B_2$};
  \node[right] at (10.28,5) {$B_i$};
  \node[right] at (10.28,9) {$B_9$};

  \node[right] at (5.35,2.95) {$P_i$};

\end{tikzpicture}
\end{center}
\nopagebreak[4]
\Needspace{0.4\textheight}
\end{question}
\begin{question}
\gkti{19}如图，在四棱柱 $ABCD-A_1B_1C_1D_1$ 中，侧棱 $AA_1\perp$ 底面 $ABCD$，$AB\parallel DC$，$AA_1=1$，$AB=3k$，$AD=4k$，$BC=5k$，$DC=6k$（$k>0$）。

（1） 求证：$CD\perp$ 平面 $ADD_1A_1$；

（2） 若直线 $AA_1$ 与平面 $AB_1C$ 所成角的正弦值为 $\dfrac67$，求 $k$ 的值；

（3） 现将与四棱柱 $ABCD-A_1B_1C_1D_1$ 形状和大小完全相同的两个四棱柱拼成一个新的四棱柱，规定：若拼成的新四棱柱形状和大小完全相同，则视为同一种拼接方案。问：共有几种不同的拼接方案？在这些拼接成的新四棱柱中，记其中最小的表面积为 $f(k)$，写出 $f(k)$ 的解析式。（直接写出答案，不必说明理由）
\nopagebreak[4]
\begin{center}% ============================================================
% f2013_fujian_science_q19_prism.tex
% 2013 福建卷（理工类）第 19 题 四棱柱（TikZ 矢量重绘）
% 侧棱 AA1 垂直底面 ABCD；AB 平行 DC
% AB=3k, AD=4k, BC=5k, DC=6k, AA1=1
% 图中另绘平面 AB1C 的三边 AB1、B1C、AC
% ============================================================
\begin{tikzpicture}[
  x=1cm, y=1cm, line width=0.6pt, font=\footnotesize, >=Stealth
]

  % ---------- 顶点（斜二测投影） ----------
  \coordinate (A)  at (0,0);
  \coordinate (B)  at (3,0);
  \coordinate (C)  at (8,1.30);
  \coordinate (D)  at (2,1.30);
  \coordinate (A1) at (0,1.15);
  \coordinate (B1) at (3,1.15);
  \coordinate (C1) at (8,2.45);
  \coordinate (D1) at (2,2.45);

  % ---------- 底面 ----------
  \draw (A) -- (B) -- (C);
  \draw[dashed] (C) -- (D) -- (A);

  % ---------- 侧棱 ----------
  \draw (A) -- (A1);
  \draw (B) -- (B1);
  \draw (C) -- (C1);
  \draw[dashed] (D) -- (D1);

  % ---------- 上底面 ----------
  \draw (A1) -- (B1) -- (C1) -- (D1) -- cycle;

  % ---------- 平面 AB1C ----------
  \draw (A) -- (B1);
  \draw (B1) -- (C);
  \draw[dashed] (A) -- (C);

  % ---------- 标注 ----------
  \node[below left]  at (A)  {$A$};
  \node[below]       at (B)  {$B$};
  \node[right]       at (C)  {$C$};
  \node[above right] at (D)  {$D$};
  \node[left]        at (A1) {$A_1$};
  \node[above]       at (B1) {$B_1$};
  \node[above right] at (C1) {$C_1$};
  \node[above]       at (D1) {$D_1$};

\end{tikzpicture}
\end{center}
\nopagebreak[4]
\end{question}
\begin{question}
\gkti{20}已知函数 $f(x)=\sin(\omega x+\varphi)$（$\omega>0$，$0<\varphi<\pi$）的周期为 $\pi$，图象的一个对称中心为 $\left(\dfrac\pi4,0\right)$。将函数 $f(x)$ 图象上所有点的横坐标伸长到原来的 $2$ 倍（纵坐标不变），再将得到的图象向右平移 $\dfrac\pi2$ 个单位长度后得到函数 $g(x)$ 的图象。

（1） 求函数 $f(x)$ 与 $g(x)$ 的解析式；

（2） 是否存在 $x_0\in\left(\dfrac\pi6,\dfrac\pi4\right)$，使得 $f(x_0)$，$g(x_0)$，$f(x_0)g(x_0)$ 按照某种顺序成等差数列？若存在，请确定 $x_0$ 的个数；若不存在，说明理由；

（3） 求实数 $a$ 与正整数 $n$，使得 $F(x)=f(x)+ag(x)$ 在 $(0,n\pi)$ 内恰有 $2013$ 个零点。

\end{question}
\begin{question}
\gkti{21}\textbf{三选二} 

以下 A、B、C 三题任选两题作答。

\textbf{【A】} 已知直线 $l:ax+y=1$ 在矩阵 $\boldsymbol A=\begin{pmatrix}1&2\\0&1\end{pmatrix}$ 对应的变换作用下变为直线 $l':x+by=1$。

（1） 求实数 $a$，$b$ 的值；

（2） 若点 $P(x_0,y_0)$ 在直线 $l$ 上，且 $\boldsymbol A\begin{pmatrix}x_0\\y_0\end{pmatrix}=\begin{pmatrix}x_0\\y_0\end{pmatrix}$，求点 $P$ 的坐标。

\textbf{【B】} 在平面直角坐标系中，以坐标原点为极点，$x$ 轴的非负半轴为极轴建立极坐标系。已知点 $A$ 的极坐标为 $\left(\sqrt2,\dfrac\pi4\right)$，直线 $l$ 的极坐标方程为 $\rho\cos\left(\theta-\dfrac\pi4\right)=a$，且点 $A$ 在直线 $l$ 上。

（1） 求 $a$ 的值及直线 $l$ 的直角坐标方程；

（2） 圆 $C$ 的参数方程为 $\begin{cases}x=1+\cos\alpha,\\y=\sin\alpha,\end{cases}$（$\alpha$ 为参数），试判断直线 $l$ 与圆 $C$ 的位置关系。

\textbf{【C】} 设不等式 $|x-2|<a$（$a\in\mathbf N^*$）的解集为 $A$，且 $\dfrac32\in A$，$\dfrac12\notin A$。

（1） 求 $a$ 的值；

（2） 求函数 $f(x)=|x+a|+|x-2|$ 的最小值。

\end{question}